\section{Архитектура системы}

\subsection{Технологический стек бэкенда}

\begin{itemize}
  \item Язык: \textbf{Python 3.13+}.
  \item Фреймворк: \textbf{FastAPI}.
  \item ОРМ: \textbf{SQLAlchemy} с поддержкой асинхронного режима.
  \item Миграции БД: \textbf{Alembic}.
  \item Валидация и конфигурация: \textbf{pydantic}, \textbf{pydantic-settings}.
  \item Драйвер БД: \textbf{asyncpg} (асинхронный доступ к PostgreSQL).
  \item Контейнеризация: \textbf{Docker}.
  \item ASGI-сервер: \textbf{Uvicorn}.
  \item Тесты: \textbf{pytest}, плагины, \textbf{testcontainers}.
  \item Хранилище файлов: \textbf{S3-совместимое} (например, MinIO, AWS S3 или корпоративное решение ИТМО).
\end{itemize}

\subsection{Фронтенд}

Фронтенд реализуется на \textbf{React} и \textbf{TypeScript}. Конкретный выбор библиотек (UI-компоненты, состояние, HTTP-клиент, просмотр PDF) фиксируется в ходе разработки и документируется (например, Axios, react-pdf или PDF.js).

\subsection{Развёртывание и обратный прокси}

Сервис развёртывается на \textbf{инфраструктуре Университета ИТМО}. Приложение работает в контейнерах Docker. Перед приложением используется \textbf{Nginx} в качестве \textbf{обратного прокси}: завершение TLS, маршрутизация запросов к бэкенду и при необходимости к статике фронтенда.

\subsection{Аутентификация}

Интеграция с \textbf{itmo.id} для входа и получения идентификатора пользователя. При недоступности itmo.id используется механизм, согласованный с инфраструктурой ИТМО.

\subsection{API}

RESTful API; описание в формате OpenAPI (Swagger). Безопасность загружаемых PDF обеспечивается проверками по п.\,3.1.4 и отражается в нефункциональных требованиях по безопасности.

\subsection{Интерфейс на основе электронной почты (компонент)}

Реализация интерфейса по п.\,3.5 может представлять собой отдельный сервис или фоновый процесс: подключение к почтовому серверу (IMAP с возможностью OAuth2 для корпоративной почты ИТМО), периодический опрос или обработка событий, фильтрация по темам писем, извлечение вложений PDF, вызов API проверки, формирование отчёта и отправка ответного письма. Конфигурация (темы писем, привязка типа документа, адрес для отчётов) задаётся через настройки приложения.

\subsection{CI/CD и автоматизированное развёртывание}

\subsubsection{Контур CI/CD}

В контуре непрерывной интеграции (например, GitLab CI или GitHub Actions) при каждом push и при создании merge request выполняются:
\begin{itemize}
  \item \textbf{Тесты:} модульные и интеграционные тесты бэкенда (pytest, testcontainers), тесты фронтенда (Jest, React Testing Library), при необходимости~--- сценарии E2E.
  \item \textbf{SAST:} статический анализ кода на уязвимости и нежелательные паттерны (например, Bandit, Semgrep или SonarQube для Python; ESLint и аналоги для TypeScript). Результаты могут приводить к падению или предупреждению пайплайна в соответствии с политикой.
\end{itemize}

\subsubsection{Развёртывание на сервере}

При успешном прохождении пайплайна для ветки main (или тега релиза) выполняется автоматическое развёртывание на целевом сервере. Развёртывание выполняется с помощью \textbf{Ansible} (playbook’и и/или роли): обновление артефактов приложения (образы Docker или собранный фронтенд/бэкенд), при необходимости переконфигурация Nginx, перезапуск сервисов, применение миграций БД. Развёртывание является частью CI/CD и не выполняется вручную в штатном режиме. Этапы пайплайна: lint, test, SAST, build, deploy. Перед приложением на сервере используется Nginx в качестве обратного прокси.

